% theme: synthesis_of_trigonometric_functions
\MajorQuestionLesson{三角関数の合成}
\SetRowSpace{4mm}

\TypeQuestion{次の式を $r\sin(\theta+\alpha)$ の形に表しなさい。ただし、$r>0$、$-\pi<\alpha\leqq\pi$ とする。}{synthesis_sine_form}
\QQRow{$\sin\theta+\cos\theta$}{}{$\sqrt{3}\sin\theta-\cos\theta$}{}
\QQRow{$-\sin\theta+\sqrt{3}\cos\theta$}{}{$\sin\theta-\cos\theta$}{}
\QQRow{$-2\sin\theta-2\cos\theta$}{}{$\cos\theta+\sqrt{3}\sin\theta$}{}
\QQRow{$2\sqrt{3}\cos\theta-2\sin\theta$}{}{$\dfrac{1}{2}\sin\theta-\dfrac{\sqrt{3}}{2}\cos\theta$}{}
\QQRow{$-\sqrt{3}\sin\theta-\cos\theta$}{}{$\dfrac{3}{2}\sin\theta+\dfrac{3}{2}\cos\theta$}{}

\TypeQuestion{次の式を $r\cos(\theta-\alpha)$ の形に表しなさい。ただし、$r>0$、$-\pi<\alpha\leqq\pi$ とする。}{synthesis_cosine_form}
\QQRow{$\sqrt{3}\cos\theta+\sin\theta$}{}{$\cos\theta-\sin\theta$}{}
\QQRow{$-\sin\theta-\sqrt{3}\cos\theta$}{}{$2\sin\theta-2\cos\theta$}{}

\TypeQuestion{次の式を $r\sin(\theta+\alpha)$ の形に表し、$r,\ \sin\alpha,\ \cos\alpha$ の値を求めなさい。ただし、$r>0$、$-\pi<\alpha\leqq\pi$ とする。}{synthesis_nonstandard_phase}
\QQRow{$3\sin\theta-4\cos\theta$}{}{$-2\sin\theta+3\cos\theta$}{}
\QQ{$\dfrac{5}{2}\sin\theta+6\cos\theta$}{}

\TypeQuestion{次の等式がすべての $\theta$ について成り立つとき、定数 $a,\ b$ の値を求めなさい。}{synthesis_determine_coefficient}
\QQ{$a\sin\theta+b\cos\theta=4\sin\left(\theta-\dfrac{\pi}{6}\right)$}{}
\QQ{$a\sin\theta+b\cos\theta=3\cos\left(\theta+\dfrac{\pi}{4}\right)$}{}

\LessonQuestionBreak

\TypeQuestion{次の関数のグラフを $0\leqq\theta\leqq2\pi$ の範囲でかき、周期を求めなさい。}{synthesis_graph}
\QQRow{$y=\sin\theta+\cos\theta$}{}{$y=\sqrt{3}\cos\theta-\sin\theta$}{}
\QQRow{$y=2\sin\theta-2\cos\theta+1$}{}{$y=\sin2\theta+\sqrt{3}\cos2\theta$}{}

\TypeQuestion{次の方程式を、$0\leqq\theta<2\pi$ の範囲で解きなさい。}{synthesis_equation}
\QQRow{$\sin\theta+\cos\theta=1$}{}{$\sqrt{3}\sin\theta-\cos\theta=1$}{}
\QQRow{$\sin\theta-\cos\theta=-\sqrt{2}$}{}{$\sin\theta+\sqrt{3}\cos\theta=\sqrt{3}$}{}
\QQRow{$\sqrt{3}\cos\theta-\sin\theta=-1$}{}{$2\sin\theta+2\cos\theta=-\sqrt{2}$}{}
\QQRow{$\sin2\theta-\sqrt{3}\cos2\theta=1$}{}{$\sqrt{3}\sin\theta+\cos\theta=-2$}{}

\TypeQuestion{次の不等式を、$0\leqq\theta<2\pi$ の範囲で解きなさい。}{synthesis_inequality}
\QQRow{$\sin\theta+\cos\theta>1$}{}{$\sqrt{3}\sin\theta-\cos\theta\leqq0$}{}
\QQRow{$\sin\theta-\cos\theta\geqq-1$}{}{$\sqrt{3}\cos\theta-\sin\theta<-1$}{}
\QQRow{$\sin\theta+\sqrt{3}\cos\theta\geqq\sqrt{3}$}{}{$\sin2\theta+\cos2\theta<1$}{}

\LessonQuestionBreak

\TypeQuestion{次の関数の最大値と最小値を、それぞれそのときの $\theta$ の値とともに求めなさい。}{synthesis_basic_extrema}
\QQRow{$y=\sin\theta+\cos\theta\quad(0\leqq\theta\leqq2\pi)$}{}{$y=\sqrt{3}\sin\theta-\cos\theta\quad(0\leqq\theta\leqq2\pi)$}{}
\QQ{$y=2\sin\theta-2\cos\theta+3\quad(0\leqq\theta\leqq\pi)$}{}
\QQ{$y=\sin\theta+\sqrt{3}\cos\theta\quad\left(0\leqq\theta\leqq\dfrac{\pi}{2}\right)$}{}
\QQ{$y=\sqrt{3}\sin\theta+\cos\theta\quad\left(-\dfrac{\pi}{6}\leqq\theta\leqq\dfrac{\pi}{6}\right)$}{}
\QQ{$y=\sin2\theta-\cos2\theta\quad\left(0\leqq\theta\leqq\dfrac{\pi}{2}\right)$}{}

\TypeQuestion{$a$ を実数とする。次の方程式が指定された範囲に解をもつような $a$ の範囲を求めなさい。}{synthesis_parameter_existence}
\QQ{$\sin\theta+\sqrt{3}\cos\theta=a\quad(0\leqq\theta<2\pi)$}{}
\QQ{$2\sin\theta-2\cos\theta=a\quad\left(0\leqq\theta\leqq\dfrac{\pi}{2}\right)$}{}
\QQ{$\sqrt{3}\sin\theta+\cos\theta=a\quad\left(0\leqq\theta\leqq\dfrac{\pi}{3}\right)$}{}
